% CSE 715 — Computer Graphics
% Chapter 5: Two-Dimensional Viewing and Clipping — Complete Model Answers (2020–2024)
% University of Chittagong, 7th Semester B.Sc. Engineering
% Source: Schaum's Outline of Computer Graphics (Xiang & Plastock), Chapter 5
%
% Compile with: pdflatex Chapter5_Solutions.tex (run twice for TOC)

\documentclass[12pt, a4paper]{article}

\usepackage[left=2.5cm, right=2.5cm, top=2.8cm, bottom=2.8cm, headheight=15pt]{geometry}
\usepackage{amsmath, amssymb}
\usepackage{tikz}
\usetikzlibrary{arrows.meta, positioning, shapes.geometric, calc, backgrounds, fit}
\usepackage{booktabs}
\usepackage{array}
\usepackage{xcolor}
\usepackage{enumitem}
\usepackage{fancyhdr}
\usepackage{titlesec}
\usepackage{mdframed}
\usepackage{multicol}
\usepackage{tabularx}
\usepackage{multirow}
\usepackage{hyperref}

%----- Colours -----
\definecolor{sectionblue}{RGB}{10,60,110}
\definecolor{boxbg}{RGB}{242,247,255}
\definecolor{answerbg}{RGB}{240,255,240}
\definecolor{notesbg}{RGB}{255,250,230}
\definecolor{keyblue}{RGB}{0,80,160}
\definecolor{accent}{RGB}{180,40,40}

%----- Box environments -----
\mdfdefinestyle{solutionframe}{linecolor=sectionblue, backgroundcolor=boxbg, roundcorner=4pt, linewidth=1.4pt, innertopmargin=7pt, innerbottommargin=7pt, innerleftmargin=9pt, innerrightmargin=9pt}
\mdfdefinestyle{answerframe}{linecolor=green!55!black, backgroundcolor=answerbg, roundcorner=4pt, linewidth=1.4pt, innertopmargin=7pt, innerbottommargin=7pt, innerleftmargin=9pt, innerrightmargin=9pt}
\mdfdefinestyle{notesframe}{linecolor=orange!70!black, backgroundcolor=notesbg, roundcorner=4pt, linewidth=1pt, innertopmargin=5pt, innerbottommargin=5pt, innerleftmargin=8pt, innerrightmargin=8pt}
\newenvironment{solutionbox}{\begin{mdframed}[style=solutionframe]}{\end{mdframed}}
\newenvironment{answerbox}{\begin{mdframed}[style=answerframe]}{\end{mdframed}}
\newenvironment{notesbox}{\begin{mdframed}[style=notesframe]}{\end{mdframed}}

%----- Page style -----
\pagestyle{fancy}
\fancyhf{}
\lhead{\small\color{sectionblue} CSE 715 --- Chapter 5: 2D Viewing \& Clipping}
\rhead{\small\color{sectionblue} Model Solutions (2020--2024)}
\cfoot{\small\thepage}
\renewcommand{\headrulewidth}{0.5pt}

%----- Section formatting -----
\titleformat{\section}{\Large\bfseries\color{sectionblue}}{}{0em}{}[\vspace{-2pt}\color{sectionblue}\rule{\linewidth}{0.8pt}]
\titleformat{\subsection}{\large\bfseries\color{sectionblue!85!black}}{}{0em}{}
\titleformat{\subsubsection}{\normalsize\bfseries\color{black}}{}{0em}{}

%----- Custom list -----
\newlist{keypoints}{itemize}{1}
\setlist[keypoints]{label=$\triangleright$, leftmargin=1.5cm, itemsep=2pt, topsep=3pt}
\newlist{examlist}{enumerate}{1}
\setlist[examlist]{label=\textbf{(\roman*)}, leftmargin=1.8cm, itemsep=3pt, topsep=4pt}

%=====================================================================
\begin{document}
%=====================================================================

\begin{center}
  {\LARGE\bfseries\color{sectionblue} CSE 715: Computer Graphics}\\[0.6em]
  {\Large\bfseries Chapter 5 --- Two-Dimensional Viewing and Clipping}\\[0.3em]
  {\large Complete Model Answers (2020--2024)}\\[0.2em]
  {\normalsize University of Chittagong \quad|\quad 7\textsuperscript{th} Semester B.Sc.\ (Engg.) Examination}\\[0.3em]
  {\small Source: Schaum's Outline of Computer Graphics (Xiang \& Plastock), Chapter 5 --- notation follows the book exactly (region-code bit order, $wx/wy$, $vx/vy$, $x_{\min}/x_{\max}/y_{\min}/y_{\max}$)}
\end{center}

\vspace{0.3em}
\noindent\rule{\linewidth}{1pt}
\vspace{0.5em}

\tableofcontents
\newpage

%=====================================================================
\section{Theoretical Framework --- Reusable Across All Years}
%=====================================================================
Chapter 5 shows up as \textbf{Q4 (Section A, every year)} for Cohen--Sutherland clipping and Sutherland--Hodgman/Weiler--Atherton polygon clipping, and inside \textbf{Q5 (Section B, every year)} for the Normalization (window-to-viewport) transformation. Convex/concave polygon identification is bundled wherever the paper happens to place it (Q4, Q6, or Q7 depending on the year) but is answered here since it is §5.4 material.

%----------------------------------------------------------------------
\subsection{Window-to-Viewport (Normalization) Transformation --- \S5.1}
%----------------------------------------------------------------------
\begin{answerbox}
A \textbf{window} is specified by four world coordinates $wx_{\min}, wx_{\max}, wy_{\min}, wy_{\max}$. A \textbf{viewport} is specified by four normalized device coordinates $vx_{\min}, vx_{\max}, vy_{\min}, vy_{\max}$. To map $(wx,wy)\to(vx,vy)$ preserving relative placement:
$$vx = \frac{vx_{\max}-vx_{\min}}{wx_{\max}-wx_{\min}}(wx-wx_{\min})+vx_{\min}, \qquad vy = \frac{vy_{\max}-vy_{\min}}{wy_{\max}-wy_{\min}}(wy-wy_{\min})+vy_{\min}$$
As a translate--scale--translate composite matrix $N$ with $s_x=\dfrac{vx_{\max}-vx_{\min}}{wx_{\max}-wx_{\min}}$, $s_y=\dfrac{vy_{\max}-vy_{\min}}{wy_{\max}-wy_{\min}}$:
$$N = \begin{pmatrix}1&0&vx_{\min}\\0&1&vy_{\min}\\0&0&1\end{pmatrix}\begin{pmatrix}s_x&0&0\\0&s_y&0\\0&0&1\end{pmatrix}\begin{pmatrix}1&0&-wx_{\min}\\0&1&-wy_{\min}\\0&0&1\end{pmatrix} = \begin{pmatrix}s_x&0&-s_x\,wx_{\min}+vx_{\min}\\0&s_y&-s_y\,wy_{\min}+vy_{\min}\\0&0&1\end{pmatrix}$$
\textbf{Workstation transformation} is the same operation a second time: NDC $\to$ physical device pixels.
\textbf{Preserving aspect ratio} (only a window given, no viewport): compute window aspect $a_w=\dfrac{\text{width}}{\text{height}}$. If $a_w\ge1$ (wide), use the full $x$-extent $[0,1]$ and $y$-extent $[0,\,1/a_w]$. If $a_w<1$ (tall), use the full $y$-extent $[0,1]$ and $x$-extent $[0,\,a_w]$. Either way $s_x=s_y$ (uniform scale $\Rightarrow$ no distortion).
\end{answerbox}

%----------------------------------------------------------------------
\subsection{Point and Line Clipping --- Cohen--Sutherland Algorithm (\S5.2--5.3)}
%----------------------------------------------------------------------
\begin{answerbox}
\textbf{Point clipping:} $(x,y)$ is inside the window iff $x_{\min}\le x\le x_{\max}$ and $y_{\min}\le y\le y_{\max}$.

\textbf{4-bit region code} (exact book order --- \textbf{TOP, BOTTOM, RIGHT, LEFT}), using $\mathrm{sign}(a)=1$ if $a>0$, else $0$:
$$\text{Bit 1 (TOP)}=\mathrm{sign}(y-y_{\max}) \quad \text{Bit 2 (BOTTOM)}=\mathrm{sign}(y_{\min}-y) \quad \text{Bit 3 (RIGHT)}=\mathrm{sign}(x-x_{\max}) \quad \text{Bit 4 (LEFT)}=\mathrm{sign}(x_{\min}-x)$$
A point with code $0000$ is inside the window.

\textbf{Three clipping categories} for a line from $(x_1,y_1)$ to $(x_2,y_2)$:
\begin{enumerate}[noitemsep]
  \item \textbf{Visible} --- both endpoint codes are $0000$.
  \item \textbf{Not visible} --- bitwise AND of the two codes $\neq 0000$ (both endpoints lie outside the \emph{same} boundary).
  \item \textbf{Clipping candidate} --- neither of the above (AND $=0000$ but not both codes $0000$).
  \end{enumerate}
\textbf{Algorithm (iterative):} pick an endpoint with a nonzero code. Push each set bit across its boundary in turn: bit TOP $\to$ intersect with $y=y_{\max}$; bit BOTTOM $\to$ $y=y_{\min}$; bit RIGHT $\to$ $x=x_{\max}$; bit LEFT $\to$ $x=x_{\min}$. Replace that endpoint with the intersection point, recompute its code, and re-classify. Repeat until the segment is Visible or Not Visible.
$$\text{Intersection with vertical boundary } x=x_i:\ y_i=y_1+m(x_i-x_1) \qquad \text{with horizontal } y=y_i:\ x_i=x_1+\frac{y_i-y_1}{m},\quad m=\frac{y_2-y_1}{x_2-x_1}$$
\end{answerbox}

%----------------------------------------------------------------------
\subsection{Polygon Clipping --- Convexity, Orientation, Left/Right Test (\S5.4)}
%----------------------------------------------------------------------
\begin{answerbox}
\textbf{Convex polygon:} the segment joining any two interior points lies entirely inside the polygon. A \textbf{concave} polygon is any polygon that is not convex (at least one interior segment exits the boundary --- typically visible as a "notch" or reflex vertex $>180^\circ$).

\textbf{Positive orientation:} a polygon $P_1,\dots,P_N$ is positively oriented if traversing the vertices in order produces a \textbf{counterclockwise} circuit --- equivalently, the left hand of someone walking along any directed edge $\overrightarrow{P_{i-1}P_i}$ points \emph{inside} the polygon.

\textbf{Left/right-of-line test:} for a directed edge $A(x_1,y_1)\to B(x_2,y_2)$ and a point $P(x,y)$:
$$C = (x_2-x_1)(y-y_1) - (y_2-y_1)(x-x_1)$$
$C>0 \Rightarrow P$ is to the \textbf{left} of $\overrightarrow{AB}$; $C<0 \Rightarrow P$ is to the \textbf{right}. For a positively-oriented convex clipping polygon: $P$ outside if to the right of \emph{any} edge, inside if to the left of \emph{every} edge.
\end{answerbox}

%----------------------------------------------------------------------
\subsection{Sutherland--Hodgman Algorithm (\S5.4)}
%----------------------------------------------------------------------
\begin{answerbox}
Clips a (possibly concave) subject polygon against one edge $E$ (endpoints $A,B$) of a positively-oriented \emph{convex} clipping polygon at a time, re-feeding the output into the next edge. For each subject edge $\overline{P_{i-1}P_i}$, using the left/right test above against $E$:
\begin{enumerate}[noitemsep]
  \item Both $P_{i-1},P_i$ \textbf{left} of $E$ $\Rightarrow$ output $P_i$.
  \item Both \textbf{right} of $E$ $\Rightarrow$ output nothing.
  \item $P_{i-1}$ left, $P_i$ right (\textbf{exiting}) $\Rightarrow$ output the intersection point $I$ of $\overline{P_{i-1}P_i}$ with (the extended) edge $E$.
  \item $P_{i-1}$ right, $P_i$ left (\textbf{entering}) $\Rightarrow$ output $I$, then $P_i$.
\end{enumerate}
Process all edges of the subject polygon (closing back to $P_1$), then repeat against the clip polygon's next edge. \textbf{Caveat:} the algorithm can join disconnected output pieces with spurious double-back edges when the subject polygon splits into multiple parts --- these coincide in opposite directions and can be removed in a post-processing pass, or avoided entirely by using Weiler--Atherton.
\end{answerbox}

%----------------------------------------------------------------------
\subsection{Weiler--Atherton Algorithm (\S5.4)}
%----------------------------------------------------------------------
\begin{answerbox}
Handles concave clip windows and correctly outputs multiple disjoint sub-polygons (Sutherland--Hodgman's weakness). Call the clip window the \textbf{clip polygon} and the polygon being clipped the \textbf{subject polygon}.
\begin{enumerate}[noitemsep]
  \item Start at any subject-polygon vertex; trace its boundary.
  \item If an edge \textbf{enters} the clip polygon, record the intersection point and continue tracing the subject polygon.
  \item If an edge \textbf{leaves} the clip polygon, record the intersection point, then make a \textbf{right turn} onto the clip polygon boundary (swap roles: now trace what was the clip polygon, treating it as the subject).
  \item Whenever the trace closes into a loop, output that closed loop as one result sub-polygon; resume from the next not-yet-traced recorded intersection.
  \item Terminate when the entire original subject-polygon border has been traced exactly once.
\end{enumerate}
\end{answerbox}

%=====================================================================
\section{Examination 2024 --- Q5(a,b,c) \& Q6(a) (Section B)}
%=====================================================================
\begin{solutionbox}
\textbf{Question breakdown:}
\begin{enumerate}[noitemsep]
  \item[Q5(a)] Find the normalization transformation that maps a window with lower-left corner $(1,1)$ and upper-right corner $(3,5)$ onto: (i) a viewport that is the entire NDC screen; (ii) a viewport with lower-left corner $(0,0)$ and upper-right corner $(v_x,v_y)$. (3)
  \item[Q5(b)] Clipping window $x_{\min}=10,x_{\max}=30,y_{\min}=10,y_{\max}=25$. Compute the 4-bit Cohen--Sutherland region code (TOP, BOTTOM, RIGHT, LEFT) for $P_1(5,5),P_2(20,30),P_3(35,15),P_4(15,8),P_5(12,20)$. Then determine whether $A(5,5)$--$B(12,20)$ is visible, not visible, or requires clipping. (3)
  \item[Q5(c)] \emph{(Related --- Circular Clipping, not core §5 topic but bundled in this question)} Circular window center $C(10,10)$, radius $r=6$; points $p_1(2,8)$, $p_2(8,16)$. Determine if the segment is completely inside, completely outside, or intersecting; if intersecting, compute the intersection point(s) and the clipped visible portion. (3)
  \item[Q6(a)] Clipping window: convex quadrilateral $A(1,1),B(5,1),C(5,4),D(1,4)$. Subject polygon: triangle $P(2,2),Q(4,2),R(3,5)$. Perform Weiler--Atherton polygon clipping: (i) show entry/exit points; (ii) list clipped-polygon vertices in CCW order. (3)
\end{enumerate}
\end{solutionbox}

\subsection*{Q5(a) --- Normalization, Window $(1,1)$--$(3,5)$ [3 marks]}
$wx_{\min}=1,wx_{\max}=3,wy_{\min}=1,wy_{\max}=5$ (width $2$, height $4$).
\begin{solutionbox}
\textbf{(i) Viewport = entire NDC} ($vx_{\min}=0,vx_{\max}=1,vy_{\min}=0,vy_{\max}=1$): $s_x=\frac{1}{2},\ s_y=\frac{1}{4}$.
$$N = \begin{pmatrix}\tfrac12&0&-\tfrac12\\0&\tfrac14&-\tfrac14\\0&0&1\end{pmatrix}$$
\textbf{(ii) Viewport $(0,0)$--$(v_x,v_y)$} (symbolic): $s_x=\dfrac{v_x}{2},\ s_y=\dfrac{v_y}{4}$.
$$N = \begin{pmatrix}\tfrac{v_x}{2}&0&-\tfrac{v_x}{2}\\0&\tfrac{v_y}{4}&-\tfrac{v_y}{4}\\0&0&1\end{pmatrix}$$
\end{solutionbox}

\subsection*{Q5(b) --- Region Codes + Visibility Classification [3 marks]}
\begin{solutionbox}
Window: $x_{\min}=10,x_{\max}=30,y_{\min}=10,y_{\max}=25$. Bit order (TOP,BOTTOM,RIGHT,LEFT):
\begin{center}
\begin{tabular}{lcccccc}
\toprule
Point & $x,y$ & TOP & BOTTOM & RIGHT & LEFT & Code \\
\midrule
$P_1$ & $(5,5)$ & 0 & 1 & 0 & 1 & \textbf{0101} \\
$P_2$ & $(20,30)$ & 1 & 0 & 0 & 0 & \textbf{1000} \\
$P_3$ & $(35,15)$ & 0 & 0 & 1 & 0 & \textbf{0010} \\
$P_4$ & $(15,8)$ & 0 & 1 & 0 & 0 & \textbf{0100} \\
$P_5$ & $(12,20)$ & 0 & 0 & 0 & 0 & \textbf{0000 (inside)} \\
\bottomrule
\end{tabular}
\end{center}
\textbf{$A(5,5)=P_1$, $B(12,20)=P_5$:} codes $0101$ AND $0000 = 0000$, and not both $0000$ $\Rightarrow$ \textbf{requires clipping} (Category 3). $A$ is outside on both LEFT and BOTTOM; clip against $x=x_{\min}=10$ first: direction $B-A=(7,15)$, $t=\frac{10-5}{7}=\frac57$, $y=5+15\cdot\frac57=\frac{110}{7}\approx15.71$. Intersection $I_1=(10,\,\tfrac{110}{7})$, code $0000$ (inside). Segment $I_1$--$B$: both $0000$ $\Rightarrow$ visible. \textbf{Visible portion: $(10,\,15.71)$ to $(12,20)$.}
\end{solutionbox}

\subsection*{Q5(c) --- Circular Clipping Window [3 marks, Related]}
\begin{solutionbox}
Circle: $(x-10)^2+(y-10)^2=36$. Check endpoints: $p_1(2,8)$: $(2-10)^2+(8-10)^2=64+4=68>36$ (outside). $p_2(8,16)$: $(8-10)^2+(16-10)^2=4+36=40>36$ (outside). Both endpoints are outside, but the chord may still pass through the circle --- parametrize $p_1\to p_2$: $x=2+6t,\ y=8+8t,\ t\in[0,1]$. Substituting:
$$(6t-8)^2+(8t-2)^2=36 \ \Rightarrow\ 100t^2-128t+32=0 \ \Rightarrow\ 25t^2-32t+8=0 \ \Rightarrow\ t=\frac{16\pm2\sqrt{14}}{25}$$
$t_1\approx0.341,\ t_2\approx0.939$ --- both in $[0,1]$, so the segment \textbf{intersects} the circle at two points (enters then exits):
$$I_1\approx(4.04,\,10.73) \qquad I_2\approx(7.64,\,15.51)$$
The clipped (visible, inside-circle) portion is the chord from $I_1$ to $I_2$; the two end stubs $p_1$--$I_1$ and $I_2$--$p_2$ lie outside the circle and are discarded.
\end{solutionbox}

\subsection*{Q6(a) --- Weiler--Atherton: Triangle vs.\ Rectangle [3 marks]}
\begin{solutionbox}
Clip window $ABCD$: $x\in[1,5],\ y\in[1,4]$. Subject triangle $P(2,2),Q(4,2),R(3,5)$: $P,Q$ inside; $R$ outside (above $y_{\max}=4$).
\textbf{Exit point} (edge $Q\to R$ leaves through top $y=4$): direction $R-Q=(-1,3)$, $t=\frac{4-2}{3}=\frac23$, $x=4-\frac23=\frac{10}{3}$. $E_1=\left(\tfrac{10}{3},4\right)$.
\textbf{Entry point} (edge $R\to P$ re-enters through top $y=4$): direction $P-R=(-1,-3)$, $t=\frac{4-5}{-3}=\frac13$, $x=3-\frac13=\frac83$. $E_2=\left(\tfrac{8}{3},4\right)$.
Trace: $P(2,2)\to Q(4,2)$ [inside] $\to$ exit at $E_1$ $\to$ right turn onto window's top edge $\to$ entry at $E_2$ $\to$ resume subject trace back to $P$.
\textbf{Clipped polygon (CCW):} $P(2,2),\ Q(4,2),\ E_1\!\left(\tfrac{10}{3},4\right),\ E_2\!\left(\tfrac{8}{3},4\right)$ --- a quadrilateral (verified positively-oriented via the shoelace formula).
\end{solutionbox}

%=====================================================================
\section{Examination 2023 --- Q4(a) (Section A) \& Q8(a) (Section B)}
%=====================================================================
\begin{solutionbox}
\textbf{Question breakdown:}
\begin{enumerate}[noitemsep]
  \item[Q4(a)] Polygon $A(1,1),B(4,1),C(4,4),D(1,4)$ (CCW). Clipping window: rectangle $W(2,2),X(3,2),Y(3,3),Z(2,3)$. Using Sutherland--Hodgman, calculate the clipped-polygon vertices step by step for each edge of the clipping window.
  \item[Q8(a)] Perform the Cohen--Sutherland algorithm for line segment endpoints $(5,5)$ and $(15,10)$ in a clipping window with boundaries $x_{\min}=0,y_{\min}=0,x_{\max}=10,y_{\max}=10$. (3.5)
\end{enumerate}
\end{solutionbox}

\subsection*{Q4(a) --- Sutherland--Hodgman: $3\times3$ Square vs.\ $1\times1$ Window}
\begin{solutionbox}
Subject: $A(1,1),B(4,1),C(4,4),D(1,4)$. Clip window edges (each processed in turn against the whole running polygon):
\textbf{Clip against left edge $x=2$} (keep $x\ge2$): $A{\to}B$ exits/enters at $(2,1)$, keep $B(4,1)$; $B{\to}C$ both in, keep $C(4,4)$; $C{\to}D$ exit at $(2,4)$; $D{\to}A$ both out. Result: $(2,1),(4,1),(4,4),(2,4)$.
\textbf{Clip against bottom edge $y=2$} (keep $y\ge2$): result $(4,2),(4,4),(2,4),(2,2)$.
\textbf{Clip against right edge $x=3$} (keep $x\le3$): result $(3,4),(2,4),(2,2),(3,2)$.
\textbf{Clip against top edge $y=3$} (keep $y\le3$): result $(2,3),(2,2),(3,2),(3,3)$.
\textbf{Final clipped polygon: $(2,3),(2,2),(3,2),(3,3)$} --- exactly the window $Z,W,X,Y$ (cyclically shifted), confirming the correct result since the $1\times1$ window lies entirely inside the $3\times3$ square.
\end{solutionbox}

\subsection*{Q8(a) --- Cohen--Sutherland: $(5,5)$--$(15,10)$, Window $(0,0)$--$(10,10)$ [3.5 marks]}
\begin{solutionbox}
$P_1(5,5)$: code $0000$ (inside). $P_2(15,10)$: $x=15>10\Rightarrow$ RIGHT$=1$; $y=10$ not $>10\Rightarrow$ TOP$=0$. Code $0010$.
AND $=0000$, one endpoint inside $\Rightarrow$ clip $P_2$ against $x=x_{\max}=10$: direction $(10,5)$, $t=\frac{10-5}{10}=0.5$, $y=5+5(0.5)=7.5$. $I=(10,7.5)$, code $0000$.
\textbf{Visible segment: $(5,5)$ to $(10,7.5)$.}
\end{solutionbox}

%=====================================================================
\section{Examination 2022 --- Q4(a,b,c) (Section A) \& Q7(b) (Section B)}
%=====================================================================
\begin{solutionbox}
\textbf{Question breakdown:}
\begin{enumerate}[noitemsep]
  \item[Q4(a)] Find region codes for $A(6,2), B(3,8), C(-1,2), D(2,4)$; window $x_{\min}=2,x_{\max}=5,y_{\min}=1,y_{\max}=9$. Clip $AB$ and $CD$ using Cohen--Sutherland. (3)
  \item[Q4(b)] Using Sutherland--Hodgman, show the clipping steps for the given polygon/window figure ($P_1$--$P_8$ vs.\ shaded clipping window). (4)
  \item[Q4(c)] Find the normalization transformation from a window (lower-left/upper-right corners as given in the paper) onto the NDC screen so that aspect ratios are preserved. (2)
  \item[Q7(b)] Draw the example diagram of the convex and concave polygon. (1)
\end{enumerate}
\end{solutionbox}

\subsection*{Q4(a) --- Region Codes + Clip $AB$, $CD$ [3 marks]}
\begin{solutionbox}
Window: $x_{\min}=2,x_{\max}=5,y_{\min}=1,y_{\max}=9$.
\begin{center}
\begin{tabular}{lccccc}
\toprule
Point & TOP & BOTTOM & RIGHT & LEFT & Code \\
\midrule
$A(6,2)$ & 0 & 0 & 1 & 0 & \textbf{0010} \\
$B(3,8)$ & 0 & 0 & 0 & 0 & \textbf{0000 (inside)} \\
$C(-1,2)$ & 0 & 0 & 0 & 1 & \textbf{0001} \\
$D(2,4)$ & 0 & 0 & 0 & 0 & \textbf{0000 (inside)} \\
\bottomrule
\end{tabular}
\end{center}
\textbf{Clip $AB$:} $A$ outside (RIGHT), clip against $x=x_{\max}=5$: direction $B-A=(-3,6)$, $t=\frac{5-6}{-3}=\frac13$, $y=2+6\cdot\frac13=4$. $I=(5,4)$, code $0000$. Segment $I$--$B$: both $0000$ $\Rightarrow$ visible. \textbf{Clipped $AB=(5,4)$ to $(3,8)$.}
\textbf{Clip $CD$:} $C$ outside (LEFT), clip against $x=x_{\min}=2$: direction $D-C=(3,2)$, $t=\frac{2-(-1)}{3}=1$ --- this reaches $D$ exactly, i.e.\ the line touches the boundary precisely at $D(2,4)$. \textbf{The entire segment $CD$ except the single point $D$ is clipped away} (the line only grazes the window at its corner-boundary $x=2$).
\end{solutionbox}

\subsection*{Q4(b) --- Sutherland--Hodgman on Figure $P_1$--$P_8$ [4 marks]}
\begin{answerbox}
No numeric coordinates are given --- the figure shows an 8-vertex subject polygon $P_1,\dots,P_8$ partly overlapping a rectangular shaded clipping window. \textbf{Method (apply directly to your own paper's figure):} process one clip-window edge at a time (left, right, top, bottom in any consistent order). For each edge $E$, walk the subject polygon's edges $\overline{P_{i-1}P_i}$ using the left/right test (\S1.3): both endpoints left of $E\Rightarrow$ output $P_i$; both right $\Rightarrow$ output nothing; left-to-right (exiting) $\Rightarrow$ output the intersection point; right-to-left (entering) $\Rightarrow$ output the intersection point then $P_i$. Re-feed the output vertex list into the next clip edge. From the figure, $P_1,P_2,P_5,P_6$ sit outside the window while $P_3,P_4,P_7,P_8$ sit on/inside it --- so the edges touching the window boundary ($P_2$--$P_3$, $P_4$--$P_5$, $P_6$--$P_7$, $P_8$--$P_1$) are exactly where new intersection vertices get generated at each clipping stage.
\end{answerbox}

\subsection*{Q4(c) --- Normalization, Aspect Ratio Preserved [2 marks]}
\begin{notesbox}
\textbf{Scan note:} the printed corners are partly illegible; read here as lower-left $(1,1)$, upper-right $(3,4)$ (width $2$, height $3$) --- the only reading consistent with a valid (non-degenerate) rectangle and with "aspect ratio preserved" phrasing (matches the book's own Solved Problem 5.4 pattern). If your copy shows different corners, the method below is unchanged --- just substitute your numbers.
\end{notesbox}
\begin{solutionbox}
$wx_{\min}=1,wx_{\max}=3,wy_{\min}=1,wy_{\max}=4$: width $2$, height $3$, $a_w=\frac23<1$ (taller than wide) $\Rightarrow$ use full $y$-extent $[0,1]$, $x$-extent $\left[0,\frac23\right]$. $s_x=\frac{2/3}{2}=\frac13,\ s_y=\frac{1}{3}=\frac13$ (equal --- aspect preserved).
$$N=\begin{pmatrix}\tfrac13&0&-\tfrac13\\0&\tfrac13&-\tfrac13\\0&0&1\end{pmatrix}$$
\end{solutionbox}

\subsection*{Q7(b) --- Convex vs.\ Concave Diagram [1 mark]}
\begin{answerbox}
\textbf{Convex:} every interior-point-to-interior-point segment stays inside (e.g.\ a triangle, rectangle, regular pentagon). \textbf{Concave:} at least one such segment exits the boundary, visible as a notch/reflex vertex ($>180^\circ$ interior angle), e.g.\ an arrow or star shape. See \S1.3 definition and Fig.\ 5-7 (Schaum's).
\end{answerbox}

%=====================================================================
\section{Examination 2021 --- Q4(a,b,c) (Section A) \& Q7(d) (Section B)}
%=====================================================================
\begin{solutionbox}
\textbf{Question breakdown:}
\begin{enumerate}[noitemsep]
  \item[Q4(a)] Region codes for $M(8,11),N(16,7),P(6,8),Q(3,5)$; window from rectangle $A(4,4),B(14,4),C(14,10),D(4,10)$. Clip $MN$ and $PQ$ using Cohen--Sutherland.
  \item[Q4(b)] Polygon $abcdefghijkl$ vs.\ viewing window $MNOP$ (figure). Using Sutherland--Hodgman, clip in order: (a) left edge $PM$, (b) bottom edge $PO$. (3.25)
  \item[Q4(c)] Identify the convex and concave polygons from three given shapes; briefly explain your selection. (1.5)
  \item[Q7(d)] Find the normalization transformation that maps a window with lower-left $(1,1)$ and upper-right $(2,2)$ onto a viewport with lower-left $(0,0)$ and upper-right $(1/2,1/2)$. (4)
\end{enumerate}
\end{solutionbox}

\subsection*{Q4(a) --- Region Codes + Clip $MN$, $PQ$}
\begin{solutionbox}
Window (from $A(4,4),B(14,4),C(14,10),D(4,10)$): $x_{\min}=4,x_{\max}=14,y_{\min}=4,y_{\max}=10$.
\begin{center}
\begin{tabular}{lccccc}
\toprule
Point & TOP & BOTTOM & RIGHT & LEFT & Code \\
\midrule
$M(8,11)$ & 1 & 0 & 0 & 0 & \textbf{1000} \\
$N(16,7)$ & 0 & 0 & 1 & 0 & \textbf{0010} \\
$P(6,8)$ & 0 & 0 & 0 & 0 & \textbf{0000 (inside)} \\
$Q(3,5)$ & 0 & 0 & 0 & 1 & \textbf{0001} \\
\bottomrule
\end{tabular}
\end{center}
\textbf{Clip $MN$:} $M$ outside (TOP), clip against $y=y_{\max}=10$: direction $N-M=(8,-4)$, $t=\frac{10-11}{-4}=0.25$, $x=8+8(0.25)=10$. $I_1=(10,10)$, code $0000$. Segment $I_1$--$N$: $N$ still RIGHT-coded, clip against $x=x_{\max}=14$: direction $(6,-3)$, $t=\frac{14-10}{6}=\frac23$, $y=10-3\cdot\frac23=8$. $I_2=(14,8)$, code $0000$. \textbf{Clipped $MN=(10,10)$ to $(14,8)$.}
\textbf{Clip $PQ$:} $Q$ outside (LEFT), clip against $x=x_{\min}=4$: direction $P-Q=(3,3)$, $t=\frac{4-3}{3}=\frac13$, $y=5+3\cdot\frac13=6$. $I=(4,6)$, code $0000$. \textbf{Clipped $PQ=(4,6)$ to $(6,8)$.}
\end{solutionbox}

\subsection*{Q4(b) --- Sutherland--Hodgman on Figure $abcdefghijkl$ vs.\ $MNOP$ [3.25 marks]}
\begin{answerbox}
Figure-only question (no numeric coordinates) --- apply the method of \S1.4/\S1.5 directly to the printed figure, processing the two named edges in the required order:
\begin{enumerate}[noitemsep]
  \item \textbf{Left clip edge $\overline{PM}$:} for each subject edge, test both endpoints' left/right status against $\overline{PM}$ (the window's left boundary). Vertices to the left survive as-is; edges crossing the boundary generate a new intersection vertex; edges entirely to the right (outside) are dropped.
  \item \textbf{Bottom clip edge $\overline{PO}$:} re-feed the output list from step (i) and repeat the same left/right test against $\overline{PO}$ (the window's bottom boundary).
\end{enumerate}
The polygon's re-entrant "V" notch (visible between the labeled vertices in the middle of the figure) means several edges cross each clip boundary twice --- track each subject edge independently rather than assuming one crossing per edge.
\end{answerbox}

\subsection*{Q4(c) --- Convex/Concave Identification [1.5 marks]}
\begin{answerbox}
Three shapes given: (a) a trapezoid --- \textbf{convex} (no notches, every interior segment stays inside); (b) a zig-zag/arrow-like polygon --- \textbf{concave} (has a reflex vertex $>180^\circ$, an interior segment across the notch exits the boundary); (c) a rounded hexagon-like shape --- \textbf{convex} (all interior angles $<180^\circ$). Justify each using the definition in \S1.3: check every vertex's interior angle, or equivalently check whether any two interior points' connecting segment leaves the shape.
\end{answerbox}

\subsection*{Q7(d) --- Normalization, Window $(1,1)$--$(2,2)$ $\to$ Viewport $(0,0)$--$(1/2,1/2)$ [4 marks]}
\begin{solutionbox}
$wx_{\min}=1,wx_{\max}=2,wy_{\min}=1,wy_{\max}=2$ (unit window); $vx_{\min}=0,vx_{\max}=\frac12,vy_{\min}=0,vy_{\max}=\frac12$. $s_x=\dfrac{1/2}{1}=\frac12,\ s_y=\frac12$.
$$N=\begin{pmatrix}\tfrac12&0&-\tfrac12\\0&\tfrac12&-\tfrac12\\0&0&1\end{pmatrix}$$
\end{solutionbox}

%=====================================================================
\section{Examination 2020 --- Q5(a,b,c) \& Q6(b,c) (Group-B)}
%=====================================================================
\begin{solutionbox}
\textbf{Question breakdown:}
\begin{enumerate}[noitemsep]
  \item[Q5(a)] Find the workstation transformation that maps the NDC screen onto a physical device with $x$ extent $0$--$199$ and $y$ extent $0$--$639$, origin at the lower-left corner. (2)
  \item[Q5(b)] Find the normalization transformation mapping window $(1,1)$--$(2,2)$ onto viewport $(0,0)$--$(1/2,1/2)$. (2.75)
  \item[Q5(c)] Briefly explain the Cohen--Sutherland algorithm. (1)
  \item[Q6(b)] Region codes for $A(6,2),B(3,8),C(-1,2),D(2,4)$; window $x_{\min}=1,x_{\max}=5,y_{\min}=1,y_{\max}=7$. Clip $AB$, $CD$. (4)
  \item[Q6(c)] Identify the convex and concave polygons from four given figures; justify your selection. (2)
\end{enumerate}
\end{solutionbox}

\subsection*{Q5(a) --- Workstation Transformation [2 marks]}
\begin{solutionbox}
Source is the entire NDC unit square ($wx_{\min}=0,wx_{\max}=1,wy_{\min}=0,wy_{\max}=1$); target device $vx_{\min}=0,vx_{\max}=199,vy_{\min}=0,vy_{\max}=639$, origin at the lower-left (so no axis flip needed). $s_x=199,\ s_y=639$.
$$N = \begin{pmatrix}199&0&0\\0&639&0\\0&0&1\end{pmatrix}$$
\end{solutionbox}

\subsection*{Q5(b) --- Normalization, Window $(1,1)$--$(2,2)$ $\to$ Viewport $(0,0)$--$(1/2,1/2)$ [2.75 marks]}
\begin{solutionbox}
Identical setup and result to 2021 Q7(d) above:
$$N=\begin{pmatrix}\tfrac12&0&-\tfrac12\\0&\tfrac12&-\tfrac12\\0&0&1\end{pmatrix}$$
\end{solutionbox}

\subsection*{Q5(c) --- Briefly Explain Cohen--Sutherland [1 mark]}
\begin{answerbox}
Assign each line endpoint a 4-bit region code (TOP/BOTTOM/RIGHT/LEFT) based on its position relative to the clipping window. If both codes are $0000$ the line is visible; if the bitwise AND of the two codes is nonzero the line is entirely outside (not visible); otherwise, repeatedly replace an outside endpoint with its intersection point against the boundary indicated by its set bit(s), recomputing the code each time, until the segment resolves to visible or not-visible.
\end{answerbox}

\subsection*{Q6(b) --- Region Codes + Clip $AB$, $CD$ [4 marks]}
\begin{solutionbox}
Window: $x_{\min}=1,x_{\max}=5,y_{\min}=1,y_{\max}=7$.
\begin{center}
\begin{tabular}{lccccc}
\toprule
Point & TOP & BOTTOM & RIGHT & LEFT & Code \\
\midrule
$A(6,2)$ & 0 & 0 & 1 & 0 & \textbf{0010} \\
$B(3,8)$ & 1 & 0 & 0 & 0 & \textbf{1000} \\
$C(-1,2)$ & 0 & 0 & 0 & 1 & \textbf{0001} \\
$D(2,4)$ & 0 & 0 & 0 & 0 & \textbf{0000 (inside)} \\
\bottomrule
\end{tabular}
\end{center}
\textbf{Clip $AB$:} $A$ outside (RIGHT), clip against $x=x_{\max}=5$: direction $B-A=(-3,6)$, $t=\frac{5-6}{-3}=\frac13$, $y=2+6\cdot\frac13=4$. $I_1=(5,4)$, code $0000$. Segment $I_1$--$B$: $B$ still TOP-coded, clip against $y=y_{\max}=7$: direction $(-2,4)$, $t=\frac{7-4}{4}=0.75$, $x=5-2(0.75)=3.5$. $I_2=(3.5,7)$, code $0000$. \textbf{Clipped $AB=(5,4)$ to $(3.5,7)$.}
\textbf{Clip $CD$:} $C$ outside (LEFT), clip against $x=x_{\min}=1$: direction $D-C=(3,2)$, $t=\frac{1-(-1)}{3}=\frac23$, $y=2+2\cdot\frac23=\frac{10}{3}$. $I=(1,\tfrac{10}{3})$, code $0000$. \textbf{Clipped $CD=(1,\,3.33)$ to $(2,4)$.}
\end{solutionbox}

\subsection*{Q6(c) --- Convex/Concave Identification [2 marks]}
\begin{answerbox}
Four shapes: (i) a trapezoid --- \textbf{convex}; (ii) a rounded/dented hexagon-like shape --- \textbf{convex} if all vertex angles $<180^\circ$, else concave (check the specific figure's dent); (iii) a zig-zag arrow shape --- \textbf{concave} (reflex vertex present); (iv) a regular pentagon --- \textbf{convex}. Justification method: for each vertex, check whether the interior angle exceeds $180^\circ$ (concave) or not (convex); equivalently, test whether any interior-to-interior segment exits the polygon (\S1.3).
\end{answerbox}

%=====================================================================
\section{Quick Summary --- Exam Patterns for Chapter 5}
%=====================================================================

\begin{center}
\renewcommand{\arraystretch}{1.3}
\begin{tabular}{p{4.7cm}|p{2.7cm}|p{6.6cm}}
\toprule
\textbf{Pattern} & \textbf{Years} & \textbf{Method to memorise} \\
\midrule
Cohen--Sutherland region codes (TOP,BOTTOM,RIGHT,LEFT) + clip 2 line segments & 2020, 2021, 2022, 2023, 2024 & Compute codes via the 4 sign tests; AND both codes to classify Visible/Not-visible/Candidate; push the outside endpoint across whichever boundary its set bit indicates, iterate \\
Normalization transformation (window $\to$ viewport, both given) & 2020, 2021, 2023(none), 2024 & $s_x,s_y$ from ratio of extents; $N=$ translate--scale--translate; watch for workstation-transform variant (NDC$\to$device) \\
Normalization with aspect ratio preserved (only window given) & 2022 & Compute $a_w=$ width/height; fill the longer NDC axis fully, scale the other axis by $a_w$ so $s_x=s_y$ \\
Sutherland--Hodgman polygon clipping (numeric or figure-based) & 2021, 2022, 2023 & Process one clip edge at a time; 4-case left/right test per subject edge; re-feed output to next edge \\
Weiler--Atherton polygon clipping (handles concave/multi-part output) & 2024 & Trace subject polygon; right-turn onto clip polygon at each exit; resume subject trace at next entry; output each closed loop \\
Convex/concave polygon identification + justification & 2020, 2021, 2022 & Check every interior angle $<180^\circ$ (convex) or find one $\ge180^\circ$ (concave, reflex vertex) \\
\bottomrule
\end{tabular}
\end{center}

\begin{notesbox}
\textbf{Exam-day strategy for Chapter 5:} region-code clipping is pure mechanical bookkeeping --- compute codes, AND them, and always clip the \emph{outside} endpoint first, one boundary at a time, recomputing the code after every intersection. For polygon clipping, default to Sutherland--Hodgman unless the question explicitly says Weiler--Atherton or the clip window is concave. The 2022 aspect-ratio-preserving normalization question is the one variant without an explicit viewport --- don't panic if only a window is given, just apply the $a_w$ rule from \S1.1.
\end{notesbox}

\end{document}
